% === ACADEMIC CV (LaTeX) ===
% Final Polished Version
\documentclass[11pt]{article}
\usepackage[margin=0.5in]{geometry}
\usepackage[hidelinks]{hyperref}
\usepackage{enumitem}
\usepackage{titlesec}
\usepackage{parskip} % For clean spacing between entries

% --- FORMATTING SETUP ---
\titleformat{\section}
  {\large\bfseries\uppercase} % Uppercase and Bold
  {}
  {0em}
  {}[\titlerule] % Horizontal line after header

\titlespacing*{\section}{0pt}{8pt}{2pt} % Adjust spacing around titles

% Consistent list formatting
\setlist[itemize,1]{label=\textbf{--}, leftmargin=*, topsep=2pt, itemsep=3pt, parsep=0pt}
\setlist[itemize,2]{label=$\circ$, leftmargin=*, topsep=2pt, itemsep=2pt, parsep=0pt}

\pagenumbering{gobble}

\begin{document}

% --- HEADER ---
\begin{center}
    {\LARGE\textbf{Muhammad Haseeb}} \\ [0.4em]
    \href{mailto:haseeb099m@gmail.com}{haseeb099m@gmail.com} $\mid$
    \href{https://github.com/ha405}{github.com/ha405} $\mid$
    \href{https://www.linkedin.com/in/muhammad-haseeb-4500b3246/}{linkedin.com/in/muhammad-haseeb}
\end{center}

% --- RESEARCH INTERESTS ---
\section{Research Interests}
AI Safety; Mechanistic Interpretability; Emergent Misalignment; Domain Generalization; Efficient ML.

% --- EDUCATION ---
\section{Education}
\textbf{Lahore University of Management Sciences (LUMS)} \hfill Sep 2022 -- May 2026 \\
B.S. Computer Science \\
\textit{Relevant Coursework:} Machine Learning, Deep Learning, Computer Vision, AI on Edge Devices, Advanced Topics in ML, LLM Systems, Linear Algebra, Probability.

% --- PUBLICATIONS ---
\section{Publications}
\begin{itemize}
    \item \textbf{Mechanistic Evidence for Preserved-but-Misaligned Representations in Non-IID FedAvg (\textit{ICML 2026 Mechanistic Interpretability Workshop})}.\\
    \textbf{Muhammad Haseeb}, Salaar Masood, Muhammad Abdullah Sohail, Mohammad Fatim Shoaib, Muhammad Tahir.
    \item \textbf{BaCP: Backbone Contrastive Pruning for Preserving Representations in Extremely Sparse Neural Networks (\textit{Preprint})}.\\
    Mohammad Haroon Khawaja, \textbf{Muhammad Haseeb}, Mohammad Fatim Shoaib, Muhammad Tahir.
\end{itemize}

% --- RESEARCH EXPERIENCE ---
\section{Research Experience}

\textbf{ML Research Assistant} \hfill Jan 2025 -- May 2026\\
\textit{\href{https://lums.edu.pk/academic-report/ar2024/centre-city.php}{Centre for Urban Informatics, Technology and Policy (CITY)}, LUMS} \\
\textit{Advisors: Dr. Muhammad Tahir, Dr. Zubair Khalid}
\begin{itemize}
    \item \textbf{Emergent Misalignment:} Investigating which circuits are associated with misaligned behavior after narrow fine-tuning, whether targeted interventions can reduce it, and whether those interventions transfer across tasks and models.
    \item \textbf{Mechanistic Analysis of Federated Learning:} (\textit{Accepted at ICML 2026 Mechanistic Interpretability Workshop})
    \begin{itemize}[label=$\circ$]
        \item Studied FedAvg through the lens of Mechanistic Interpretability to analyze performance loss under Non-IID conditions using class-specific circuit discovery, Universal Sparse Autoencoders (USAEs), and linear probes.
        \item Compared class-specific circuits and observed drift and interference patterns during non-IID aggregation, while IID circuits behaved in a perfect constructive manner.
        \item Demonstrated via USAE concept analysis that useful latent features, similar to the IID feature basis, remain present and recoverable even when global accuracy degrades under label skew.
    \end{itemize}
    \item \textbf{Backbone Contrastive Pruning (BaCP):}
        \begin{itemize}[label=$\circ$]
            \item Contributed to a generalized pruning framework to merge standard pruning criteria with contrastive learning
            to prevent representational collapse.
            \item Designed a multi-objective loss function that aligns sparse embeddings with pretrained, fine-tuned, and
            historical snapshot references to maintain feature consistency with pretrained references to prevent representational collapse.
            \item Maintained accuracy comparable to dense baselines at up to \textbf{99\% sparsity}, outperforming standard unstructured pruning approaches across ViTs/CNNs/Language Models.
        \end{itemize}
    \item \textbf{Domain Generalization \& Interpretability:}
            \begin{itemize}[label=$\circ$]
                \item Applied circuit extraction techniques to isolate task-specific subnetworks, analyzing whether out-of-distribution failures are caused by shortcut learning or weight interference.
                \item Optimized visual queries in Vision Transformers using reinforcement learning (Group Relative Query Optimization), achieving a \textbf{3\% accuracy gain} on the PACS dataset compared to empirical risk minimization.
            \end{itemize}
    \item \textbf{Quantization of Diffusion Models:} Proposed Log-SNR conditioned temporal dynamic quantization (\textbf{55\% lower LPIPS} on PixArt-$\alpha$) and implemented Hessian-based GPTQ/QroNos for \textbf{4-bit weight compression}.
\end{itemize}

\textbf{AI/SWE Research Intern} \hfill May 2025 -- Aug 2025\\
\textit{University of Illinois Urbana-Champaign (UIUC)}
\begin{itemize}
    \item Investigated LLM-based heuristics for automated \texttt{\#ifdef} guard insertion in C codebases and explored large language models for software debloating and code dependency analysis.
    \item Contributed to the development of a VS Code extension supporting a C-to-Rust translation pipeline.
\end{itemize}

\textbf{ML Research Assistant} \hfill Jan 2025 -- Aug 2025\\
\textit{Computer Vision \& Graphics Lab, LUMS}
\begin{itemize}
    \item \textbf{KL Aware Quantization (KLAWQ):} Proposed an augmented GPTQ framework integrating KL divergence and second-order curvature to optimize a combined MSE+KL objective, reducing LLM perplexity by \textbf{30\%} compared to GPTQ baselines.
    \item \textbf{Single-Image Camera Calibration:} Developed a pipeline to train a Multi-Scale Deformable Transformer by combining loss with geometric constraints to recover full camera projection matrices from a single image.
\end{itemize}

% --- INDUSTRY EXPERIENCE ---
\section{Industry Experience}
\textbf{Founding Product Engineer} \hfill June 2026 -- Present\\
\textit{NEXA (Remote)}
\begin{itemize}
    \item Architecting and building highly scalable, distributed backend systems from the ground up.
    \item Integrating machine learning models and optimizing distributed serving pipelines for fast, efficient inference.
\end{itemize}

\textbf{OSS Engineer} \hfill Dec 2025 -- Feb 2026\\
\textit{DatacurveAI (YC S24)}
\begin{itemize}
    \item Solved bugs and added new features across various open-source machine learning repositories to test and train LLMs on working with large codebases.
\end{itemize}

\textbf{Machine Learning Engineer} \hfill Jul 2025 -- Oct 2025\\
\textit{Innova Tech}
\begin{itemize}
  \item Automated data annotation pipeline, refined training configurations, and fixed model architectures getting a \textbf{4\% accuracy gain}. Also, implemented inference optimization with TensorRT for speedup and various quantization techniques to reduce memory footprint by \textbf{40\%} on edge devices.
\end{itemize}

% --- TEACHING EXPERIENCE ---
\section{Teaching Experience}
\textbf{Teaching Fellow, MSAI} \hfill Aug 2026 -- Present\\
\textit{AI 630: Responsible AI Engineering, LUMS}
\begin{itemize}
    \item Developed and graded quizzes for AI 630: Responsible AI Engineering.
    \item Managed class logistics and coordinated course projects.
\end{itemize}

\textbf{Teaching Assistant} \hfill Aug 2025 -- Dec 2025\\
\textit{CS436: Computer Vision, LUMS}
\begin{itemize}
    \item Designed and supervised an assignment on transfer learning and multithreaded C++ object detection. Mentored 40 student groups in building a virtual tour application using Structure from Motion.
\end{itemize}

% --- SKILLS & CONTRIBUTIONS ---
\section{Skills \& Open-Source}
\textbf{Technical Skills:} Python, C, C++, Rust, SQL; PyTorch, Transformers, ONNX, TensorRT, diffusers, adapters.

\textbf{Open-Source Contributions:}
\begin{itemize}[label=$\circ$, leftmargin=1.5em]
    \item \textbf{pytorch-image-models:} Added F1, precision, and recall metrics to evaluation and training pipelines.
    \item \textbf{adapters:} Implemented PEFT support for Group Query Attention models, fixing tensor mismatch issues.
\end{itemize}

\end{document}
